% Part: first-order-logic
% Chapter: introduction
% Section: semantic-notions

\documentclass[../../../include/open-logic-section]{subfiles}

\begin{document}

\olfileid{fol}{int}{sem}

\olsection{语义概念}

我们之前提到，当考虑$\Sat{M}{!A}[s]$是否成立时，我们（为了方便）让 $s$ 为所有变元指派值，但实际用到的只是它指派给 $!A$ 中变元的值。事实上，只有 $!A$ 中\emph{自由}变元的值才真正起作用。当然，出于严谨，我们接下来将证明这一事实。由于语句没有自由变元，在考虑它们是否在一个结构中被满足时，$s$ 就完全无关紧要了。因此，当 $!A$ 是语句时，我们可以定义 $\Sat{M}{!A}$ 意指“$\Sat{M}{!A}[s]$对所有 $s$ 成立”，而这实际上当且仅当存在至少一个 $s$ 使得 $\Sat{M}{!A}[s]$ 成立。我们需要引入变元指派以得到适用于公式的满足关系的定义，但对语句而言，满足关系独立于变元指派。

一旦我们有了“$\Sat{M}{!A}$”的定义，就一阶逻辑的语句而言，我们就知道了“情形”和“在……中为真”是什么意思。基于对语句 $\Sat{M}{!A}$ 的定义，我们可以进而定义有效性、语义后承和可满足性这些基本的语义概念。如果一个语句被所有结构满足，则称该语句是有效的，记作 $\Entails !A$。如果每个满足 $\Gamma$ 中所有语句的结构也都满足 $!A$，则称一组语句 $\Gamma$ 语义后承 $!A$，记作 $\Gamma \Entails !A$。如果存在某个结构同时满足其中的所有语句，则称一组语句是可满足的。

因为公式是归纳定义的，而满足关系又是基于公式的结构归纳定义的，所以我们可以使用归纳法来证明我们语义的性质，并建立所定义的这些语义概念之间的联系。我们将收集并证明其中一些性质，部分原因是它们本身具有理论趣味，但主要是因为在我们后续研究语义与推导系统的关系时，其中许多性质将发挥重要作用。为此，我们还必须（精确地，即通过归纳法）定义一些尚未提及的句法概念和操作。

\end{document}
